\documentclass[oneside,a4paper,14pt]{extarticle}
\usepackage[T2A]{fontenc}
\usepackage[utf8]{inputenc}
\usepackage[russian]{babel}
\usepackage[a4paper,top=2cm,bottom=2cm,left=2.5cm,right=1.5cm]{geometry}
\usepackage{textcomp}
\usepackage{indentfirst}
\usepackage{listings}
\usepackage{graphicx}
\usepackage{mwe}
\usepackage{wrapfig}
\usepackage{caption}
\usepackage{amsmath}
\usepackage{amsfonts}
\usepackage{amsthm}
\usepackage[all]{xy}
\usepackage[breaklinks]{hyperref}
\usepackage{xurl}

%% Настройка заголовков разделов
\usepackage{titlesec}
\titleformat{\section}{\normalsize\bfseries}{\thesection}{1em}{}
\titleformat{\subsection}{\normalsize\bfseries}{\thesubsection}{1em}{}
\titleformat{\subsubsection}{\normalsize\bfseries}{\thesubsubsection}{1em}{}

\renewcommand{\baselinestretch}{1.33}
\normalsize
\setlength{\parindent}{1.25cm}

\begin{document}
	\newpage\thispagestyle{empty}
	\begin{center}
		МИНИСТЕРСТВО НАУКИ И ВЫСШЕГО ОБРАЗОВАНИЯ\\
		РОССИЙСКОЙ ФЕДЕРАЦИИ
		ФЕДЕРАЛЬНОЕ ГОСУДАРСТВЕННОЕ БЮДЖЕТНОЕ\\
		ОБРАЗОВАТЕЛЬНОЕ
		УЧРЕЖДЕНИЕ ВЫСШЕГО ОБРАЗОВАНИЯ\\
		«ВЯТСКИЙ ГОСУДАРСТВЕННЫЙ УНИВЕРСИТЕТ»\\
		Институт математики и информационных систем\\
		Факультет автоматики и вычислительной техники\\
		Кафедра электронных вычислительных машин
	\end{center}
	\vspace{20mm}

	\begin{center}
		Реферат\\
		по дисциплине\\
		<<Физическая культура>>\\
		Тема: \\
		<<Статические упражнения в развитии студента>>\\
	\end{center}
	\vspace{48mm}

	\noindent
	Выполнил студент гр. ИВТб-1303-06-00 \hfill \rule[-0,5mm]{30mm}{0.15mm}\,/Гортоломей И.К./
	\newline
  
	\vfill
	\begin{center}
		Киров\\
		2026
	\end{center}
% \newpage\thispagestyle{empty}
% \begin{center}
% МИНИСТЕРСТВО НАУКИ И ВЫСШЕГО ОБРАЗОВАНИЯ\\
% РОССИЙСКОЙ ФЕДЕРАЦИИ\\[1ex]
% ФЕДЕРАЛЬНОЕ ГОСУДАРСТВЕННОЕ БЮДЖЕТНОЕ\\
% ОБРАЗОВАТЕЛЬНОЕ УЧРЕЖДЕНИЕ ВЫСШЕГО ОБРАЗОВАНИЯ\\
% <<ВЯТСКИЙ ГОСУДАРСТВЕННЫЙ УНИВЕРСИТЕТ>>\\[1ex]
% Институт математики и информационных систем\\
% Факультет автоматики и вычислительной техники\\
% Кафедра электронных вычислительных машин
% \end{center}

% \vspace{15mm}
% \begin{center}
% \Large Реферат\\[1ex]
% по дисциплине\\
% <<Физическая культура>>\\[1ex]
% Тема:\\
% <<Статические упражнения в развитии студента>>
% \end{center}

% \vspace{35mm}
% \noindent Выполнил студент гр. ИВТб-1303-06-00 \hfill \rule{4cm}{0.4pt}\,/\,Гортоломей~И.К./

% \vfill
% \begin{center}
% Киров\\
% 2025
% \end{center}

\newpage

% ==================== ВВЕДЕНИЕ ====================
\section*{ВВЕДЕНИЕ}

На современном этапе развития системы высшего образования вопросы сохранения и укрепления здоровья студенческой молодёжи приобретают особую остроту и значимость. В условиях стремительно меняющегося информационного общества, возросших академических нагрузок, а также доминирования малоподвижного образа жизни, связанного с длительным пребыванием за компьютером и в аудиториях, физическое состояние студентов претерпевает заметные изменения. В связи с этим перед педагогическим сообществом встаёт закономерный вопрос о поиске наиболее эффективных, доступных и физиологически обоснованных средств физического воспитания, способных гармонично вписываться в напряжённый учебный график и при этом обеспечивать комплексное развитие личности. В данном контексте статические (изометрические) упражнения представляют собой один из наиболее перспективных, но при этом недостаточно полно раскрытых в массовой студенческой практике инструментов.

Актуальность темы обусловлена тем, что традиционные динамические формы физической активности не всегда могут быть полноценно реализованы в условиях ограниченного времени, отсутствия специализированного оборудования или индивидуальных медицинских противопоказаний. Статические нагрузки, напротив, отличаются высокой адаптивностью, минимальными требованиями к пространству и выраженной направленностью на развитие силовой выносливости, стабилизации опорно-двигательного аппарата и формирования волевых качеств. Кроме того, в современной  научно-педа-\\гогической литературе отмечается устойчивая тенденция к переосмыслению роли физического воспитания в вузе: из узкоприкладной дисциплины оно трансформируется в важнейший компонент профессионально-личностного становления будущего специалиста. В этой связи изучение влияния статических упражнений на физическое, психоэмоциональное и профессионально-прикладное развитие студента приобретает не только теоретическую, но и практическую значимость.

Целью данной работы является комплексное рассмотрение роли и места статических упражнений в системе физического воспитания студентов высших учебных заведений, а также выявление их влияния на физиологические, психологические и профессионально-ориентированные параметры развития обучающегося.

Для достижения поставленной цели были сформулированы следующие задачи:
\begin{enumerate}
    \item Раскрыть сущность и физиологические механизмы воздействия статических упражнений на организм человека.
    \item Проанализировать влияние изометрических нагрузок на психоэмоциональную сферу и волевые качества студентов.
    \item Изучить методические особенности применения статических упражнений в учебно-воспитательном процессе вуза.
    \item Сформулировать практические рекомендации по интеграции статических нагрузок в самостоятельную и аудиторную физкультурную деятельность студентов.
\end{enumerate}

Объектом исследования выступает процесс физического воспитания студентов в условиях высшего учебного заведения. Предметом исследования являются статические упражнения как средство комплексного развития студенческой молодёжи.

В работе использованы методы теоретического анализа научно-методической литературы, обобщения педагогического опыта, а также систематизации и классификации данных. Структура реферата обусловлена логикой исследования и включает введение, три главы, заключение и список использованных источников.

% ==================== ГЛАВА 1 ====================
\section{ГЛАВА. ТЕОРЕТИЧЕСКИЕ АСПЕКТЫ СТАТИЧЕСКИХ УПРАЖНЕНИЙ В СИСТЕМЕ ФИЗИЧЕСКОГО ВОСПИТАНИЯ}

\subsection{Понятийный аппарат и классификация статических нагрузок}

В современной теории физической культуры и спорта под статическими (изометрическими) упражнениями принято понимать форму мышечной деятельности, при которой напряжение мышечных волокон происходит без видимого изменения длины мышцы и без движения в суставах. В отличие от динамических сокращений, сопровождающихся концентрической или эксцентрической фазой, статический режим характеризуется фиксацией тела или его отдельных звеньев в определённой позе на протяжении заданного временного интервала. Следует отметить, что подобный тип нагрузки издавна присутствует в различных системах физического совершенствования: от классической йоги и гимнастики ушу до современных методик функционального тренинга и реабилитационной физкультуры. Тем не менее, именно в контексте студенческого физического воспитания изометрические режимы долгое время оставались на периферии внимания, уступая место традиционным беговым, игровым и силовым динамическим комплексам.

Классификация статических упражнений может осуществляться по различным основаниям. По степени вовлечённости мышечных групп выделяют локальные (напряжение отдельных мышц или небольших мышечных цепей), региональные (нагрузка на крупные мышечные массивы одной анатомической зоны) и глобальные (удержание позы, требующей одновременной работы большинства скелетных мышц). По длительности удержания нагрузки принято различать кратковременные (до 15--20 секунд), средние (20--40 секунд) и продолжительные (свыше 40--60 секунд) статические режимы. Кроме того, в педагогической практике активно применяются упражнения со статодинамическим компонентом, где фазы неподвижного удержания чередуются с минимальными амплитудными движениями, что позволяет сочетать преимущества обоих режимов мышечной работы.

\subsection{Физиологические механизмы адаптации к изометрическим нагрузкам}

С точки зрения физиологии, статические упражнения запускают в организме ряд специфических адаптационных реакций, которые принципиально отличаются от реакций на динамическую работу. При изометрическом напряжении происходит значительное сдавление внутримышечных сосудов, что ведёт к временному ограничению кровотока в работающих мышечных группах. Данный факт обусловливает переход мышечного метаболизма на анаэробный путь энергоснабжения, сопровождающийся накоплением лактата и стимуляцией гипертрофии мышечных волокон, в первую очередь быстрых (гликолитических). При этом важно подчеркнуть, что энергозатраты при статическом удержании позы распределяются неравномерно: значительная часть энергии расходуется на поддержание нейромышечной активации и стабилизацию суставов, а не на преодоление внешнего сопротивления, что делает данный вид нагрузки относительно экономичным в калорийном эквиваленте, но высокоэффективным с точки зрения развития локальной силовой выносливости.

Нельзя не отметить и влияние статических упражнений на сердечно-сосудистую и дыхательную системы. При длительном мышечном напряжении наблюдается рефлекторное повышение артериального давления и частоты сердечных сокращений, что связано с активацией симпатического отдела вегетативной нервной системы и механическим затруднением венозного возврата. Тем не менее, при грамотном дозировании нагрузки и соблюдении принципов правильного дыхания (непрерывный выдох в фазе напряжения, отсутствие задержек) данные реакции носят тренирующий характер и способствуют повышению резервных возможностей организма. Многочисленные исследования в области спортивной физиологии подтверждают, что систематическое применение изометрических режимов способствует укреплению сухожильно-связочного аппарата, улучшению нейромышечной проводимости и формированию устойчивых проприоцептивных связей, что особенно актуально для студентов, чья профессиональная деятельность впоследствии будет сопряжена с длительными статическими позами или монотонными двигательными стереотипами.

Таким образом, физиологическая база статических упражнений представляет собой сложный, но хорошо изученный механизм адаптации, который при методически грамотном применении способен обеспечить значительный тренировочный эффект без необходимости использования громоздкого оборудования или значительных временных затрат.

% ==================== ГЛАВА 2 ====================
\section{ГЛАВА. ВЛИЯНИЕ СТАТИЧЕСКИХ НАГРУЗОК НА ПСИХОФИЗИОЛОГИЧЕСКОЕ И ПРОФЕССИОНАЛЬНОЕ РАЗВИТИЕ СТУДЕНТА}

\subsection{Формирование волевых качеств и стрессоустойчивости}

Особое значение в контексте студенческого развития имеет психологический аспект выполнения статических упражнений. В отличие от динамических нагрузок, где утомление наступает постепенно и часто маскируется ритмичностью движений, статическое удержание позы сопровождается острым, локализованным ощущением жжения, дрожи и нарастающего дискомфорта. Преодоление данного состояния требует от студента осознанного волевого усилия, способности к саморегуляции и концентрации внимания на внутреннем состоянии. В педагогической практике этот феномен нередко описывается как <<тренировка терпения в миниатюре>>, поскольку каждый повтор статического упражнения представляет собой микро-ситуацию принятия решения: продолжить удержание или прекратить выполнение. Регулярное прохождение подобных микро-стрессов способствует формированию устойчивой мотивационно-волевой сферы, развитию дисциплинированности и повышению порога переносимости психоэмоциональных нагрузок.

Следует особо подчеркнуть, что в условиях учебной деятельности, характеризующейся сессиями, дедлайнами, защитой проектов и высоким уровнем информационной перегрузки, способность студента сохранять самообладание и управлять своим внутренним состоянием приобретает ключевое значение. Статические упражнения, по своей сути являющиеся практикой осознанного контроля над телом и дыханием, косвенно тренируют те же нейронные контуры, которые задействуются при стрессоустойчивости, целеполагании и эмоциональной регуляции. Многие исследователи в области педагогической психологии справедливо указывают на то, что студенты, систематически практикующие изометрические нагрузки, демонстрируют более высокие показатели академической настойчивости, реже склонны к прокрастинации и эффективнее справляются с ситуационным тревожным фоном.

\subsection{Профессионально-прикладная направленность и профилактика учебных рисков}

Важным аспектом является профессионально-прикладная составляющая статического тренинга. Современный рынок труда предъявляет высокие требования к физической и психической устойчивости выпускников, причём для большинства специальностей характерны либо длительные статические позы (офисные работники, программисты, дизайнеры, бухгалтеры), либо эпизодические, но интенсивные статодинамические нагрузки (строители, медики, логисты, инженеры-технологи). В этой связи физическое воспитание в вузе не может ограничиваться лишь общеукрепляющей функцией; оно должно готовить организм студента к специфике будущей профессиональной деятельности. Статические упражнения, направленные на укрепление мышечного корсета, мышц спины, шеи, плечевого пояса и нижних конечностей, выполняют важнейшую профилактическую роль, минимизируя риски развития профессионально обусловленных заболеваний: остеохондроза, синдрома запястного канала, хронической мышечной усталости, нарушений осанки и венозного застоя.

Кроме того, нельзя игнорировать и тот факт, что статический тренинг обладает выраженной компенсаторной функцией. Длительное пребывание в сидячем положении, характерное для лекционных занятий и самостоятельной работы, приводит к укорочению подвздошно-поясничных мышц, ослаблению ягодичных и глубоких стабилизаторов позвоночника, а также к нарушению биомеханики дыхания. Включение в распорядок дня даже коротких (10--15 минут) комплексов изометрических упражнений позволяет частично нивелировать данные негативные эффекты, улучшить микроциркуляцию в проблемных зонах и восстановить мышечный баланс. В результате студент не только сохраняет работоспособность, но и формирует устойчивую привычку к осознанному отношению к своему телу, что в долгосрочной перспективе способствует сохранению трудоспособности и качества жизни на протяжении всей профессиональной карьеры.

Таким образом, влияние статических нагрузок выходит далеко за рамки узкофизкультурного контекста, интегрируясь в систему личностного, психологического и профессионального становления будущего специалиста.

% ==================== ГЛАВА 3 ====================
\section{ГЛАВА. МЕТОДИЧЕСКИЕ ОСОБЕННОСТИ И ПРАКТИЧЕСКАЯ РЕАЛИЗАЦИЯ В УЧЕБНОМ ПРОЦЕССЕ ВУЗА}

\subsection{Принципы дозирования и безопасности}

Эффективность применения статических упражнений в студенческой среде напрямую зависит от соблюдения научно обоснованных методических принципов. В первую очередь речь идёт о принципе постепенности и последовательности. Организм студента, особенно в первые месяцы обучения или после длительного перерыва в занятиях, не готов к резкому увеличению объёма изометрической нагрузки. Рекомендуется начинать с удержания базовых поз (планка, <<стульчик>>, мостик, статический присед) в течение 10--15 секунд с 2--3 повторами, постепенно увеличивая время фиксации на 5--10\% еженедельно. Подобный подход позволяет минимизировать риски микротравм, избежать выраженного крепатурного синдрома и сформировать устойчивый нейромышечный паттерн.

Не менее важным является принцип контроля дыхания. Как уже отмечалось ранее, задержка дыхания при статическом напряжении ведёт к резкому скачку внутригрудного и внутрибрюшного давления, что может спровоцировать нежелательные реакции со стороны сердечно-сосудистой системы, особенно у лиц со скрытыми патологиями или недостаточным уровнем тренированности. Поэтому методически грамотное руководство должно настаивать на непрерывном, ровном дыхании с акцентом на удлинение выдоха в фазе максимального напряжения. Кроме того, обязательным условием является предварительная разминка и последующая растяжка работающих мышечных групп, что обеспечивает профилактику мышечных спазмов и способствует более быстрому восстановлению.

\subsection{Интеграция в учебный процесс и самостоятельную деятельность}

В практической плоскости статические упражнения могут быть интегрированы в различные формы физкультурной деятельности студентов. Во-первых, они целесообразно вписываются в вводную часть занятий по физической культуре, выполняя функцию активации глубоких мышц-стабилизаторов перед основной динамической работой. Во-вторых, изометрические комплексы могут использоваться в качестве самостоятельных микро-тренировок в перерывах между лекциями или во время подготовки к экзаменам, что особенно актуально в условиях плотного расписания. В-третьих, статические элементы активно применяются в реабилитационных группах здоровья, где динамические нагрузки противопоказаны или ограничены по медицинским показаниям.

Опыт ряда передовых вузов свидетельствует о том, что наибольший педагогический эффект достигается при комбинированном подходе, где статические упражнения сочетаются с элементами дыхательной гимнастики, миофасциального релиза и когнитивных задач на концентрацию. Подобная интеграция не только повышает мотивацию студентов, но и формирует у них целостное представление о здоровье как о динамическом балансе физического, психического и социального благополучия. Важно также отметить, что преподаватели физического воспитания должны осуществлять регулярный мониторинг самочувствия обучающихся, корректировать нагрузку в зависимости от индивидуальной подготовленности и своевременно вносить изменения в программу при появлении признаков перетренированности или снижения адаптационных резервов.

Таким образом, методически выверенное внедрение статических упражнений в учебный процесс вуза представляет собой не просто техническое добавление новых движений, а стратегический шаг в сторону персонализированного, здоровьесберегающего и профессионально-ориентированного физического воспитания.

% ==================== ЗАКЛЮЧЕНИЕ ====================
\section*{ЗАКЛЮЧЕНИЕ}

Подводя итоги проведённого исследования, можно с уверенностью констатировать, что статические упражнения представляют собой ценный, научно обоснованный и педагогически оправданный инструмент в системе физического воспитания студенческой молодёжи. Их физиологическое воздействие, основанное на изометрическом режиме мышечной работы, обеспечивает эффективное развитие силовой выносливости, укрепление мышечного корсета, улучшение проприоцепции и стабилизацию опорно-двигательного аппарата без необходимости использования специализированного инвентаря или значительных временных затрат. В то же время психологический компонент статического тренинга способствует формированию волевых качеств, повышению стрессоустойчивости, развитию навыков самоконтроля и эмоциональной регуляции, что непосредственно коррелирует с успешностью учебной и будущей профессиональной деятельности.

Кроме того, практическая реализация статических нагрузок в вузовской среде демонстрирует высокую адаптивность: они могут быть легко интегрированы в аудиторные занятия, самостоятельные тренировки, реабилитационные программы и профессионально-прикладную подготовку. При соблюдении базовых методических принципов (постепенности, систематичности, контроля дыхания, индивидуализации и комплексности) изометрические упражнения становятся безопасным и высокоэффективным средством гармонизации физического и психического состояния студента.

Вместе с тем, следует признать, что потенциал статических упражнений в современном высшем образовании раскрыт ещё не в полной мере. Дальнейшие исследования в данной области целесообразно направить на изучение долгосрочных эффектов комбинированных статодинамических программ, разработку цифровых инструментов самоконтроля для студентов, а также на адаптацию методик под специфику различных профессиональных направлений. Только комплексный, научно-педагогический и индивидуально-ориентированный подход позволит превратить физическое воспитание из формальной учебной дисциплины в действенный механизм формирования здоровой, устойчивой и профессионально успешной личности.

Таким образом, тема статических упражнений в развитии студента остаётся актуальной, многогранной и перспективной как для теоретического осмысления, так и для практического применения в современной образовательной среде.

% ==================== СПИСОК ИСТОЧНИКОВ ====================
\begin{thebibliography}{9}
\bibitem{1} Виленский, М. Я. Физическая культура и здоровый образ жизни студента : учебное пособие / М. Я. Виленский, А. Г. Горшков, В. И. Ильинич. -- Москва : КноРус, 2021. -- 312 с.
\bibitem{2} Евсеев, Ю. И. Физическая культура : учебник для вузов / Ю. И. Евсеев. -- 5-е изд., перераб. -- Москва : Спорт, 2020. -- 456 с.
\bibitem{3} Ильинич, В. И. Профессионально-прикладная физическая подготовка студентов вузов / В. И. Ильинич. -- Москва : Высшая школа, 2019. -- 288 с.
\bibitem{4} Козлов, А. А. Изометрические тренировки в системе физического воспитания молодёжи : монография / А. А. Козлов. -- Санкт-Петербург : Лань, 2022. -- 174 с.
\bibitem{5} Матвеев, Л. П. Теория и методика физической культуры : учебник / Л. П. Матвеев. -- Москва : Физкультура и спорт, 2021. -- 544 с.
\bibitem{6} Пономарёв, В. В. Психофизиологические аспекты статической выносливости у студентов / В. В. Пономарёв // Теория и практика физической культуры. -- 2023. -- № 4. -- С. 45--49.
\bibitem{7} Холодов, Ж. К. Практикум по физиологии спорта и физического воспитания / Ж. К. Холодов, В. С. Кузнецов. -- Москва : Академия, 2020. -- 384 с.
\bibitem{8} Электронный ресурс : КиберЛенинка. Режим доступа: \url{https://cyberleninka.ru} (дата обращения: 15.05.2025).
\end{thebibliography}

\end{document}